% Generated draft: requires mathematical and source review.
\documentclass[11pt]{article}
\usepackage[margin=25mm]{geometry}
\usepackage{fontspec}
\setmainfont{Latin Modern Roman}
\usepackage{amsmath,amssymb,mathtools,unicode-math}
\setmathfont{Latin Modern Math}
\usepackage{enumitem,needspace,xurl,hyperref}
\setlength{\parindent}{0pt}
\setlength{\parskip}{5pt}
\setlength{\emergencystretch}{3em}
\newcommand{\N}{\mathbb{N}}
\newcommand{\Z}{\mathbb{Z}}
\newcommand{\Q}{\mathbb{Q}}
\newcommand{\R}{\mathbb{R}}
\newcommand{\C}{\mathbb{C}}
\newcommand{\F}{\mathbb{F}}
\newcommand{\A}{\mathbb{A}}
\newcommand{\PP}{\mathbb{P}}
\newcommand{\E}{\mathbb{E}}
\newcommand{\eps}{\varepsilon}
\newcommand{\dd}{\,\mathrm{d}}
\begin{document}
\section*{UnsolvedMath ID: 20002893}
\textbf{Source problem code:} catalogue entry
\subsection{Definitions and mathematical statement}
Let $\Omega\subsetneq\C\mathrm P^n$ be a domain with nonempty smooth boundary, pseudoconvex in the original problem. Use a fixed smooth Hermitian metric, for example the Fubini--Study metric, to define form norms and Sobolev spaces. For $0\leq p\leq n$ and $1\leq q\leq n$, Mei-Chi Shaw asks whether every
\[
f\in C^\infty_{(p,q)}(\overline\Omega),\qquad\overline\partial f=0,
\]
admits a $(p,q-1)$-form $u\in W^1(\Omega)$ satisfying $\overline\partial u=f$ in distributions. The barred operator, smoothness up to the boundary, and the exponent $\eps$ below are reconstructed from the supplied OCR-verification information; the canonical extraction loses bars and replaces $\eps$ with a control character.

The original comment says solutions are known in $L^2(\Omega)$ and in $W^\eps(\Omega)$ for some $\eps>0$ depending on $\Omega$, citing [CSW], and says the $W^1$ question was not known even with real-analytic boundary. These are supplied historical assertions, not a new assessment of the literature.

The aid reports an unconditional top-degree case: for every proper smoothly bounded domain $\Omega\subset\C\mathrm P^n$ and every $0\leq p\leq n$, there is $C=C(\Omega,p)>0$ such that every $f\in L^2_{(p,n)}(\Omega)$ has a primitive $u\in W^1_{(p,n-1)}(\Omega)$ with
\[
\overline\partial u=f,\qquad\|u\|_{W^1}\leq C\|f\|_{L^2}.
\]
Pseudoconvexity and a separate closedness hypothesis on top-degree $f$ are unnecessary for this asserted endpoint result. The presence of a genuine boundary is retained; the statement is not extended to the entire compact projective space. The aid leaves the general pseudoconvex problem open in degrees $1\leq q\leq n-1$.

Its finite-order reduction concerns the smooth-data assertion exactly. Define the closed Fr\'echet space
\[
Z^\infty_{p,q}(\overline\Omega)=\{f\in C^\infty_{(p,q)}(\overline\Omega):\overline\partial f=0\}
\]
with the standard $C^m$ seminorms up to the boundary. For each fixed $\Omega,p,q$, the assertion
\[
\text{every }f\in Z^\infty_{p,q}(\overline\Omega)\text{ has a }W^1\text{ primitive}
\]
is reported equivalent to the existence of an integer $m\geq0$ and $C>0$ such that for every such $f$ one can choose a primitive satisfying
\[
\|u\|_{W^1}\leq C\|f\|_{C^m(\overline\Omega)}.
\]
The finite-order norm may equivalently be a maximum of finitely many defining seminorms. Constants and $m$ may depend on the domain and bidegree. This is an existence estimate from an open-mapping argument, not an assertion of a continuous linear right inverse, and it does not yield an $L^2$-to-$W^1$ estimate in the unresolved degrees.
\subsection{Short English statement}
Ask for a one-derivative Sobolev solution of $\overline\partial u=f$ for smooth closed data on a pseudoconvex projective domain. The supplied aid isolates an unrestricted top-degree estimate and reformulates the remaining qualitative assertion as a finite-order smooth-data existence bound.
\subsection{Sources}
\begin{itemize}
\item Local UnsolvedMath registry, permanent ID $20002893$, source code \texttt{catalogue entry}; canonical \texttt{aim-several-complex-variables-notes.json}, \texttt{notes[50]}.
\item AIM workshop \emph{Complexity of mappings in CR geometry}, Problem 11, \emph{Solutions to $\overline\partial$}, proposed by Mei-Chi Shaw: \url{https://aimath.org/WWN/complexitycrmap/complexitycrmap.pdf}. The source cites [CSW] for known $L^2$ and $W^\eps$ solvability; no complete entry is supplied. Remarks and original literature field are empty.
\item The supplied verification records missing operator/closure bars, flattened superscripts and OCR control character U+000F in place of epsilon. Statement status is \texttt{reconstructed\_unverified}.
\item Research attempt 1: \texttt{partial\_result}, status \texttt{partially\_solved}, classification \texttt{PARTIAL-PROGRESS}, rigor \texttt{proved}, novelty confidence low. The endpoint estimate, open-mapping reduction and literature assessment are machine-generated aid, not peer reviewed, and require independent verification before citation.
\end{itemize}
\end{document}
